% !TeX root = ../papers/nc-p3.tex
\section{四次中心の持ち上げと閉点の有限判定}
\label{sec:klein-quartic-center-lift}
第二補完の既約族のうち、上側の中心行列の対角を1に規格化できる開集合を取る。
節\ref{sec:klein-doublet-boundary}の記号で
\[
 z=\frac{1-c}{b-1},\quad Z=\begin{pmatrix}1&b\\c&1\end{pmatrix},
 \quad S=Z+\operatorname{diag}(v,w),\quad
 S'=Z+\frac{1-bc}{vw}\operatorname{diag}(v,w)
\]
である。$U$は全成分1で、$V$は右下の$z$以外が1である。
分母と二つの厚いブロックの行列式は非零とする。

\begin{theorem}\label{thm:klein-quartic-center-slice}
$v=4,w=6$の二次元断面では一般の四次中心の文字次数別次元は$(3,2,0,0)$である。
$(b,c)=(2,3)$を含む単一のザリスキー開集合上、全ての閉点が標準Koszul左右AS正則で、
各始点の代数は四つの中心四次式の多項式環上で階数256の自由加群となる。
\end{theorem}
\begin{proof}
厚い二次数について、各六関係の基底を先の組への射影が単位行列となるよう選ぶ。
$X=\left(\begin{smallmatrix}1&1\\1&z\end{smallmatrix}\right)$,
$Y=\left(\begin{smallmatrix}1&1\\z&1\end{smallmatrix}\right)$とすると、後の組への
二射影は第一窓で$YS^{-1},(S'^{-1}X)^{\mathsf t}$、相手窓で
$(S^{-1}X)^{\mathsf t},Y(S')^{-1}$である。これは
式\eqref{eq:doublet-completions}の列空間・行空間を取ったものである。
各始点の次数4と5の商基底を作り、内部次数4の元$q_p$に
$q_p x_i=x_i q_{p+1}$を課す。\path{computations/verify_klein_center_lift.py}は
$\mathbb Q(b,c)$上で関係を挿入して厳密に掃き出す。文字次数0では36列の階数33、
文字次数1では40列の階数38である。上記の有理点では全四次数の核次元が
$(3,2,0,0)$で、各始点の内部次数5までの商次元も標準値である。
第二補完の全ての特殊化で標準Hilbert関数が下界なので、次数5までの商はこの点の
近傍で次数別・文字別にベクトル束である。従って他の二文字次数の核は一般に零である。
同じ近傍を縮めれば、五つの中心元の空間は次数別の部分束である。

この有理点で文字次数0から二式、文字次数1から二式を取った商の次元は
定理\ref{thm:klein-rational-period-four}と同じ素数1009の計算で
$(1+t+t^2+t^3)^4$の係数以下で、次数13で零になる。
定階数の核束を用いて四式を局所切断に持ち上げる。商の次数13までの次元の上界は
有限個の行列の階数下界なので共通の開近傍で保たれる。中心性による多項式環上の
自由加群からの全射と標準Hilbert関数の下界を比べれば全次数で一致する。
同定理の双対完全性の証明は同じ四頂点でそのまま適用できる。
\end{proof}

\begin{proposition}\label{prop:klein-central-strata}
第二補完の閉包の任意の入力で、隣接階数6と一文字階数4を仮定する。
四つの中心四次式で上記の自由加群表示を得る点の集合は、有限個の階数計算で
判定できる構成可能集合であり、その全点で標準AS正則となる。
\end{proposition}
\begin{proof}
内部次数5までの商の階数を層別化し、各定階数小行列の上で四次の中心方程式の
階数を計算する。各核の基底の一次結合を四つ、係数を独立の不定元として作る。
四式を加えた商の次元を次数13まで計算し、各始点で上の母関数の係数以下かを判定する。
これは係数について開条件であり、無限体上で一組が存在することと有理関数体上で
成立することは同値である。基底や斉次な文字次数の選択を追加制約とする必要はない。
成立すれば前定理と同じ上下界が使える。逆に標準Hilbert級数を持つ中心四変数上の
自由表示があれば、その変数を割った母関数は必ず上の多項式である。
この条件が成立しない入力のAS性は、この判定だけからは従わない。
\end{proof}
